% =====================================================================
%  FICHE 4 — THEME 2 — Corrige niveau MOYEN
% =====================================================================
\documentclass[11pt,a4paper]{article}
\usepackage[utf8]{inputenc}
\usepackage[T1]{fontenc}
\usepackage{lmodern}
\usepackage[french]{babel}
\usepackage{amsmath,amssymb,mathtools}
\usepackage{icomma}
\usepackage[version=4]{mhchem}
\usepackage{siunitx}
\sisetup{output-decimal-marker={,},per-mode=symbol}
\DeclareSIUnit\Molar{M}
\usepackage{xcolor}
\usepackage{tikz}
\usepackage[most]{tcolorbox}
\usepackage{enumitem}
\usepackage{booktabs}
\usepackage{array}
\usepackage{fancyhdr}
\usepackage[a4paper,margin=2.2cm,headheight=26pt,headsep=10pt]{geometry}

\definecolor{primary}{RGB}{146,95,160}
\definecolor{primarydark}{RGB}{92,52,104}
\definecolor{excol}{RGB}{0,135,90}

\tcbset{boxbase/.style={enhanced,breakable,boxrule=0.9pt,arc=2.5pt,
   left=3mm,right=3mm,top=2mm,bottom=2mm,fonttitle=\bfseries,coltitle=white,
   toptitle=0.5mm,bottomtitle=0.5mm}}
\newcounter{exo}
\newenvironment{corr}[1][]{%
  \par\medskip\refstepcounter{exo}%
  \noindent\begin{tcolorbox}[boxbase,colback=excol!5,colframe=excol,
     title={\textsc{Correction \theexo}\ifx\relax#1\relax\relax\else\textmd{ --- \itshape #1}\fi}]}%
  {\end{tcolorbox}}

\pagestyle{fancy}\fancyhf{}
\fancyhead[L]{\small\textcolor{primarydark}{\textbf{Fiche 4} --- Th.\,2 : Masse volumique, densit\'e, gaz}}
\fancyhead[R]{\small\textcolor{primarydark}{\textsc{Chimie} --- \textbf{Corrigé Moyen}}}
\fancyfoot[C]{\small\thepage}
\renewcommand{\headrulewidth}{0.8pt}
\renewcommand{\headrule}{\hbox to\headwidth{\color{primary}\leaders\hrule height \headrulewidth\hfill}}
\newcommand{\NA}{N_{\!\text{A}}}
\setlength{\parindent}{0pt}

\begin{document}

\begin{center}
{\LARGE\bfseries\color{primarydark} Thème 2 --- Corrigé Moyen}
\end{center}

\begin{corr}[chaîne volume $\to$ masse $\to$ quantité]
\textbf{a)} Densité $1{,}70\Leftrightarrow\rho=\SI{1.70}{\gram\per\milli\litre}$, donc
$m=\rho\,V=1{,}70\times100=\SI{170}{\gram}$.\\[2pt]
\textbf{b)} $M(\ce{CH2Cl2})=12+2(1)+2(35{,}5)=12+2+71=\SI{85}{\gram\per\mole}$, d'où
\[ n=\frac{m}{M}=\frac{170}{85}=\SI{2.0}{\mole}. \]
\[ \boxed{n(\ce{CH2Cl2})=\SI{2.0}{\mole}} \]
\end{corr}

\begin{corr}[masse molaire d'un gaz]
\textbf{a)} Pour un gaz, $n=\dfrac{V}{V_m}$ :
\[ n=\frac{5{,}6}{22{,}4}=\SI{0.25}{\mole}. \]
\textbf{b)} La masse molaire est la masse d'une mole, soit $M=\dfrac{m}{n}$ :
\[ M=\frac{5{,}0}{0{,}25}=\SI{20}{\gram\per\mole}. \]
\[ \boxed{M=\SI{20}{\gram\per\mole}} \]
\end{corr}

\begin{corr}[masse volumique et conversions « piégeuses »]
On convertit d'abord la masse : $\SI{7.72}{\hecto\gram}=7{,}72\times100=\SI{772}{\gram}$.
On isole le volume dans $\rho=m/V$ :
\[ V=\frac{m}{\rho}=\frac{772}{19{,}3}=\SI{40}{\cubic\centi\metre}=\SI{40}{\milli\litre}. \]
\[ \boxed{V=\SI{40}{\milli\litre}} \]
\textit{Piège :} oublier de convertir le hectogramme ($\SI{1}{\hecto\gram}=\SI{100}{\gram}$) mène à
un résultat 100 fois trop petit.
\end{corr}

\begin{corr}[du gaz au volume]
$M(\ce{O2})=\SI{32}{\gram\per\mole}$. On passe par la quantité de matière :
\[ n=\frac{16}{32}=\SI{0.50}{\mole}\ \Rightarrow\ V=n\,V_m=0{,}50\times22{,}4=\SI{11.2}{\litre}. \]
\[ \boxed{V=\SI{11.2}{\litre}} \]
\end{corr}

\begin{corr}[masse volumique par déplacement d'eau]
\textbf{a)} Le volume du solide est l'\textbf{augmentation} du niveau d'eau :
$V=15{,}0-10{,}0=\SI{5.0}{\milli\litre}$.\\[2pt]
\textbf{b)} $\rho=\dfrac{m}{V}=\dfrac{20{,}0}{5{,}0}=\SI{4.0}{\gram\per\milli\litre}
=\SI{4.0}{\gram\per\cubic\centi\metre}$.
\[ \boxed{V=\SI{5.0}{\milli\litre};\quad \rho=\SI{4.0}{\gram\per\cubic\centi\metre}} \]
\textit{Piège :} prendre le niveau final ($\SI{15}{\milli\litre}$) au lieu de l'augmentation
($\SI{5}{\milli\litre}$).
\end{corr}

\end{document}
